%
% teil1.tex -- Beispiel-File für das Paper
%
% (c) 2020 Prof Dr Andreas Müller, Hochschule Rapperswil
%
% !TEX root = ../../buch.tex
% !TEX encoding = UTF-8
%


\section{Riemannscher Abbildungssatz
\label{elektro:section:ReimannSatz}}
\kopfrechts{Riemannscher Abbildungssatz}

Georg Fridrich Bernhard Riemann (1826-1866) \textcolor{red}{cite from bibliography} war ein deutscher Mathematiker, der sich mit diversen Gebieten der Mathematik beschäftigte.
Auf vielen hat er bahnbrechende Ergebnisse erzielt, die bis heute von grosser Bedeutung sind, auch Einstein hat von seinen Theorien Gebrauch gemacht.
Der Riemannsche Abbildungssatz hat er erstmals als Skizze im Jahre 1851 in seiner Dissertation veröffentlicht. 
Erst später wurde der Satz von anderen Mathematikern formal bewiesen und in die Mathematik eingeführt. 
Es wird verschiedene Beweise für den Satz gefunden, die alle auf unterschiedlichen mathematischen Gebieten basieren.
Hier wird der verbreitetste Beweis vorgestellt, der auf dem Satz von Montel basiert, der die Existenz von konformen Abbildungen auf einfach zusammenhängende Gebiete beschreibt.
Beschriebt ist hier wörtlich gemeint, es gibt für den Riemannschen Abbildungssatz keine allgemeine Formel, die für alle Gebiete gilt, sondern nur einen Existenzbeweis einer Abbildung.
Dies bedeutet, dass es für jegliche Figuren eine Transformation existiert diese muss aber auf anderem Wege ermittelt werden dazu mehr in Abschnitt \ref{elektro:section:KreisTrafo}.

Im Allgemeinen ist unser Ziel ein Gebiet $G$ auf ein einfacheres Gebiet $G'$ abzubilden, um die Berechnung von Potential Feldern und elektrischen Feldern zu vereinfachen. 
Das ziel Gebiet $G'$ ist für uns der Einheitskreis, da für diesen die Formeln für das Potentialfeld $\varphi$ und das elektrische Feld $\mathbf{E}$ bereits bekannt sind. \textcolor{red}{Abschnitt verweis}

\subsection{Voraussetzungen
\label{elektro:subsection:RiemannVoraussetzungen}}
Zunächst müssen einpaar Voraussetzungen für den Riemannschen Abbildungssatz erfüllt sein, damit der Satz überhaupt angewendet werden kann.

\subsubsection{Einfach zusammenhängende Gebiete}
Das ursprüngliche Gebiet $G$ muss einfach zusammenhängend sein, das bedeutet, dass es keine Löcher in dem Gebiet gibt. 
Das Gebiet $G$ darf beliebig gross und beliebig geformt sein, solange es keine ausgeschnittene Bereiche enthält.
Zu dieser Definition gehört auch dazu, dass das Gebiet offen sein muss. Soll heissen der Rand von dem Gebiet gehört nicht zum Gebiet dazu.
Der Rand müsste separat betrachtet werden, wobei sich meistens herausstellt, dass der Rand sich brav mit verformt und die Abbildung auf dem Rand auch konform ist.
Es entstehe erst Abweichungen, wenn der Rand ausgefallene Fraktale bildet, dies ist bei uns nicht der Fall daher wird stets angenommen das der Rand analog zur Fläche verformt wird.

\subsubsection{Teilgebiete der komplexen Ebene}
$ G \subset \mathbb{C} $ muss gelten, dies bedeutet, dass $G$ ein Subset und nicht äquivalent zu $\mathbb{C}$ sein darf.
Für die komplette komplexe Ebene gilt der Riemannsche Abbildungssatz.
Wenn die gesamte komplexe Ebene auf den Einheitskreis abgebildet werden soll, dann müssten punkte die bereits im Innern des Einheitskreises liegen, unberührt bleiben und die, die ausserhalb liegen müssten auf bereits besetzte Punkte transformeirt werden. 
Die nächste Voraussetzung würde damit verletzt werden, die Bijektivität der Abbildung.

\subsubsection{Bijektive Abbildung}
Bijektiv bedeutet, dass alles eindeutig abgebildet wird, es gibt keine Punkte, die auf den gleichen Punkt abgebildet werden. 
Es geht also sowohl bei der Transformation sowie bei der Rücktransformation nichts an Information verloren.
Ansonsten würde eine Zweideutigkeit entstehen und die Abbildung wäre nicht mehr umkehrbar.

